% ============================================================
%  cover.tex — 分析学笔记封面（九层设计）
%  主色调：BlueMain 深蓝 + Gold 暖金
%  用法：\coverimage{figure/分析封面.jpg} → \makecover
% ============================================================

% ── 封面图片命令 ──
\makeatletter
\providecommand{\@coverimage}{}
\newcommand{\coverimage}[1]{\renewcommand{\@coverimage}{#1}}
\newcommand{\makecover}{%
\begin{titlepage}
\thispagestyle{empty}

% ============================================================
%  第一层：左侧深蓝色竖条 + 渐变纹理 + 装饰圆点
% ============================================================
\begin{tikzpicture}[remember picture,overlay]
  % 主体竖条
  \fill[BlueMain!90]
    (current page.south west) rectangle
    ([xshift=2.8cm]current page.north west);
  % 顶部向下的渐变淡出
  \shade[top color=BlueDark!90, bottom color=BlueMain!90]
    ([xshift=0cm,yshift=-5cm]current page.north west) rectangle
    ([xshift=2.8cm]current page.north west);
  % 底部向上的渐变淡出
  \shade[bottom color=BlueDark!90, top color=BlueMain!90]
    ([xshift=0cm,yshift=5cm]current page.south west) rectangle
    ([xshift=2.8cm]current page.south west);
  % 竖条右边界线
  \draw[BlueLight!40,line width=1pt]
    ([xshift=2.8cm]current page.south west) --
    ([xshift=2.8cm]current page.north west);
  % 左侧装饰细线（距左边缘 0.4cm）
  \draw[BlueLight!20,line width=0.4pt]
    ([xshift=0.4cm]current page.south west) --
    ([xshift=0.4cm]current page.north west);

  % 蓝色区域内的装饰圆点
  \foreach \y/\r/\o in {
    0.08/1.1/0.06, 0.22/0.7/0.10, 0.40/1.4/0.05,
    0.55/0.55/0.09, 0.68/0.9/0.07, 0.82/0.6/0.10, 0.93/0.75/0.08
  } {
    \fill[white, opacity=\o]
      ([xshift=1.4cm, yshift=-\y\paperheight]current page.north west)
      circle (\r cm);
  }

  % 两条细竖线装饰
  \draw[white, line width=0.3pt, opacity=0.15]
    ([xshift=0.5cm]current page.north west) --
    ([xshift=0.5cm]current page.south west);
  \draw[white, line width=0.3pt, opacity=0.15]
    ([xshift=2.1cm]current page.north west) --
    ([xshift=2.1cm]current page.south west);
\end{tikzpicture}

% ============================================================
%  第二层：顶部/底部装饰金线带 + 等距标记圆点
% ============================================================
\begin{tikzpicture}[remember picture,overlay]
  % 顶部横贯金线
  \draw[GoldMain!45,line width=0.6pt]
    ([xshift=3.0cm,yshift=-1.6cm]current page.north west) --
    ([xshift=-2.0cm,yshift=-1.6cm]current page.north east);
  % 金线上方小标记点
  \foreach \x in {3.5,7,10.5,14} {
    \fill[GoldMain!55]
      ([xshift=\x cm,yshift=-1.45cm]current page.north west)
      circle (1.5pt);
  }
  % 底部横贯金线
  \draw[GoldMain!35,line width=0.5pt]
    ([xshift=3.0cm,yshift=2.2cm]current page.south west) --
    ([xshift=-1.5cm,yshift=2.2cm]current page.south east);
  % 底部金线标记
  \foreach \x in {4,8,12} {
    \fill[GoldMain!40]
      ([xshift=\x cm,yshift=2.35cm]current page.south west)
      circle (1.2pt);
  }
\end{tikzpicture}

% ============================================================
%  第三层：右上角装饰几何图形组（同心圆环 + 放射线）
% ============================================================
\begin{tikzpicture}[remember picture,overlay]
  % 大圆环（空心）
  \draw[GoldMain!30,line width=0.7pt]
    ([xshift=-2.5cm,yshift=-2.5cm]current page.north east)
    circle (22pt);
  % 小圆环（同心）
  \draw[BlueLight!25,line width=0.5pt]
    ([xshift=-2.5cm,yshift=-2.5cm]current page.north east)
    circle (14pt);
  % 中心实心点
  \fill[GoldMain!50]
    ([xshift=-2.5cm,yshift=-2.5cm]current page.north east)
    circle (2.5pt);
  % 正交放射线
  \foreach \ang in {0,90,180,270} {
    \draw[GoldMain!25,line width=0.4pt]
      ([xshift=-2.5cm,yshift=-2.5cm]current page.north east)
      -- ++(\ang:28pt);
  }
  % 对角放射线
  \foreach \ang in {45,135,225,315} {
    \draw[BlueLight!18,line width=0.3pt]
      ([xshift=-2.5cm,yshift=-2.5cm]current page.north east)
      -- ++(\ang:26pt);
  }
\end{tikzpicture}

% ============================================================
%  第四层：左下角装饰 —— 数学曲线簇（正弦型 + 指数型）
% ============================================================
\begin{tikzpicture}[remember picture,overlay]
  % 波动曲线簇
  \draw[GoldMain!15,line width=0.5pt,smooth]
    ([xshift=3.5cm,yshift=3.0cm]current page.south west)
    to[out=30,in=210] ++(3.0,2.0);
  \draw[BlueLight!12,line width=0.5pt,smooth]
    ([xshift=3.3cm,yshift=3.0cm]current page.south west)
    to[out=35,in=205] ++(3.2,2.2);
  \draw[GoldMain!12,line width=0.5pt,smooth]
    ([xshift=3.7cm,yshift=2.8cm]current page.south west)
    to[out=25,in=215] ++(2.8,1.8);
  % 渐近下降曲线
  \draw[BlueMain!10,line width=0.5pt,smooth]
    ([xshift=3.5cm,yshift=4.0cm]current page.south west)
    to[out=-15,in=195] ++(2.5,-1.5);
  \draw[GoldMain!10,line width=0.4pt,smooth]
    ([xshift=3.7cm,yshift=4.2cm]current page.south west)
    to[out=-20,in=190] ++(2.3,-1.3);
  % 流方向小箭头
  \draw[GoldMain!22,-{Stealth[length=3pt,width=3pt]}]
    ([xshift=5.0cm,yshift=4.2cm]current page.south west)
    -- ++(0.8,0.5);
  \draw[BlueLight!18,-{Stealth[length=3pt,width=3pt]}]
    ([xshift=5.2cm,yshift=3.5cm]current page.south west)
    -- ++(0.7,-0.5);
\end{tikzpicture}

% ============================================================
%  第五层：公式背景区 —— 淡金色矩形 + 分析学经典公式
% ============================================================
\begin{tikzpicture}[remember picture,overlay]
  % 主背景矩形（淡金渐变）
  \shade[top color=GoldPale!40, bottom color=GoldPale!22]
    ([xshift=3.3cm,yshift=-3cm]current page.north west) rectangle
    ([xshift=-1.5cm,yshift=2.5cm]current page.south east);
  % 矩形的细金边框
  \draw[GoldMain!28,line width=0.4pt]
    ([xshift=3.3cm,yshift=-3cm]current page.north west) rectangle
    ([xshift=-1.5cm,yshift=2.5cm]current page.south east);

  % 公式间递进连接箭头
  \draw[GoldMain!30,line width=0.4pt,-{Stealth[length=3.5pt,width=3.5pt]}]
    ([xshift=5.5cm,yshift=-4.5cm]current page.north west) --
    ([xshift=5.5cm,yshift=-5.0cm]current page.north west);
  \draw[GoldMain!30,line width=0.4pt,-{Stealth[length=3.5pt,width=3.5pt]}]
    ([xshift=5.7cm,yshift=-5.5cm]current page.north west) --
    ([xshift=5.7cm,yshift=-6.0cm]current page.north west);

  % 六条分析学经典公式
  \node[anchor=west,BlueMain!65] at ([xshift=3.7cm,yshift=-3.5cm]current page.north west)
    {$\displaystyle \mu\Bigl(\bigcup_{k=1}^{\infty}E_k\Bigr)
      =\sum_{k=1}^{\infty}\mu(E_k)\quad\text{for disjoint }(E_k)$};

  \node[anchor=west,BlueLight!55] at ([xshift=3.8cm,yshift=-4.5cm]current page.north west)
    {$\displaystyle \int_X f_n\,\dd\mu\to\int_X f\,\dd\mu,
      \quad f_n\to f\ \text{in }L^1(X)$};

  \node[anchor=west,GoldMain!70] at ([xshift=3.9cm,yshift=-5.5cm]current page.north west)
    {$\displaystyle L^p(\Omega)
      =\Bigl\{u:\Omega\to\mathbb{R}\;\Big|\;\int_{\Omega}|u|^p<\infty\Bigr\}$};

  \node[anchor=west,BlueMain!50] at ([xshift=3.7cm,yshift=-6.5cm]current page.north west)
    {$\displaystyle \|u\|_{H_0^1(\Omega)}^2
      =\int_{\Omega}|\nabla u|^2\,\dd x$};

  \node[anchor=west,BlueLight!45] at ([xshift=3.8cm,yshift=-7.5cm]current page.north west)
    {$\displaystyle \langle Tu,v\rangle_H
      =\langle u,T^\ast v\rangle_H,\quad T:H\to H\ \text{bounded}$};

  \node[anchor=west,GoldMain!55] at ([xshift=3.9cm,yshift=-8.5cm]current page.north west)
    {$\displaystyle \widehat{f}(\xi)=\int_{\mathbb{R}^n}f(x)\,e^{-2\pi i x\cdot\xi}\,\dd x$};

  % 公式编号装饰（右侧小圆点）
  \foreach \y/\c in {
    -3.5/BlueMain!30, -4.5/BlueLight!28, -5.5/GoldMain!32,
    -6.5/BlueMain!25, -7.5/BlueLight!22, -8.5/GoldMain!28
  } {
    \fill[\c]
      ([xshift=-2.0cm,yshift=\y cm]current page.north west)
      circle (1.8pt);
  }
\end{tikzpicture}

% ============================================================
%  第六层：标题区 —— 装饰分隔线 + 标题文字
% ============================================================
\vspace*{2cm}

\hspace*{3.3cm}
\begin{minipage}{0.7\textwidth}
  % ── 系列标识 ──
  {\large\bfseries\color{BlueLight} Huginn Notes in Mathematics\par}

  \vspace{0.3cm}
  % ── 装饰短金线（带渐小圆点尾饰）──
  \begin{tikzpicture}
    \draw[GoldMain!60,line width=1.2pt] (0,0) -- (2.0,0);
    \fill[GoldMain!70] (2.3,0) circle (2pt);
    \fill[GoldMain!50] (2.8,0) circle (1.2pt);
    \fill[GoldMain!30] (3.2,0) circle (0.8pt);
  \end{tikzpicture}

  \vspace{0.8cm}

  % ── 主标题 ──
  {\Huge\bfseries\color{BlueDark} 分析学笔记\par}

  \vspace{0.6cm}
  % ── 作者 ──
  {\Large\color{BlueMain!70} Notes of Huginn Wong\par}

  \vspace{0.4cm}
  % ── 第二道装饰金线（两端圆点）──
  \begin{tikzpicture}
    \draw[GoldMain!40,line width=0.7pt] (0,0) -- (3.5,0);
    \fill[GoldMain!45] (0,0) circle (1.5pt);
    \fill[GoldMain!45] (3.5,0) circle (1.5pt);
  \end{tikzpicture}

  \vspace{0.8cm}
  % ── Fraktur 铭文 ──
  {\large\color{GoldMain!80} $\mathfrak{For}\quad\mathfrak{the}\quad\mathfrak{Worthy}$\par}

  \vspace{1.0cm}

%  % ── 封面图片区域（可选插入小图）──
%  \ifdefempty{\@coverimage}{
%    % 无图片：什么都不显示
%  }{
%    \begin{center}
%      \includegraphics[width=0.55\textwidth, height=3.5cm, keepaspectratio]{\@coverimage}
%    \end{center}
%    \vspace{0.3cm}
%  }
%
%  \vfill
%
%  % ── 日期 ──
%  {\large\color{BlueLight!65} \@date\par}
%  \vspace{0.8cm}
\end{minipage}

% ============================================================
%  第七层：右下角「出版社」小 logo（仿 Springer 风格）
% ============================================================
\begin{tikzpicture}[remember picture, overlay]
  \begin{scope}[shift={([xshift=-5.5cm, yshift=1.5cm]current page.south east)}]
    % 小房子
    \draw[BlueMain, line width=0.8pt] (0,0) rectangle (0.9,0.72);
    \draw[BlueMain, line width=0.8pt] (0,0.72) -- (0.45,1.08) -- (0.9,0.72);
    \draw[GoldMain, line width=0.5pt] (0.45,0.72) -- (0.45,1.08);
    \draw[BlueMain, line width=0.4pt] (0.23,0) -- (0.23,0.72);
    \draw[BlueMain, line width=0.4pt] (0.67,0) -- (0.67,0.72);
    % 屋顶金点
    \fill[GoldMain] (0.45,1.08) circle (1.5pt);
    \node[anchor=west, BlueMain, font=\footnotesize] at (1.1,0.32)
      {Produced by Huginn};
  \end{scope}
\end{tikzpicture}
% ============================================================
%  第八层：底部主封面图片 —— 与金色框左右对齐，位于小房子上方
% ============================================================
\ifdefempty{\@coverimage}{}{
	\begin{tikzpicture}[remember picture,overlay]
		\node[
		anchor=south west,
		inner sep=0pt,
		outer sep=0pt
		] at ([xshift=3.3cm,yshift=3.05cm]current page.south west)
		{
			\includegraphics[
			width=\dimexpr\paperwidth-4.8cm\relax
			]{\@coverimage}
		};
	\end{tikzpicture}
}


% ============================================================
%  第九层：散落装饰点（页面随机分布的数学点缀）
% ============================================================
\begin{tikzpicture}[remember picture,overlay]
  % 右中区域散点
  \fill[GoldMain!18] ([xshift=-3.0cm,yshift=-9.5cm]current page.north west) circle (1.0pt);
  \fill[BlueLight!18] ([xshift=-3.8cm,yshift=-9.8cm]current page.north west) circle (0.7pt);
  \fill[GoldMain!16] ([xshift=-2.5cm,yshift=-10.2cm]current page.north west) circle (0.8pt);
  % 右下方散点
  \fill[BlueMain!12] ([xshift=-4.0cm,yshift=3.5cm]current page.south west) circle (1.2pt);
  \fill[GoldMain!13] ([xshift=-5.5cm,yshift=4.0cm]current page.south west) circle (0.9pt);
  % 标题区附近轻微装饰
  \fill[GoldMain!20] ([xshift=14.5cm,yshift=-6.0cm]current page.north west) circle (1.0pt);
  \fill[BlueLight!16] ([xshift=14.0cm,yshift=-5.5cm]current page.north west) circle (0.6pt);
  \fill[GoldMain!18] ([xshift=15.0cm,yshift=-5.8cm]current page.north west) circle (0.8pt);
\end{tikzpicture}

\end{titlepage}
}
\makeatother
